% =====================================================================
%  CARNET D'ENTRAINEMENT — FICHE 6 — THÈME 5
%  Niveau DIFFICILE — Corrigé
% =====================================================================
\documentclass[11pt,a4paper]{article}
\usepackage[utf8]{inputenc}
\usepackage[T1]{fontenc}
\usepackage{lmodern}
\usepackage[french]{babel}
\usepackage{amsmath,amssymb,mathtools}
\usepackage{icomma}
\usepackage[version=4]{mhchem}
\usepackage{siunitx}
\sisetup{output-decimal-marker={,}}
\usepackage{xcolor}
\usepackage{colortbl}
\usepackage{tikz}
\usepackage[most]{tcolorbox}
\usepackage{enumitem}
\usepackage{array}
\usepackage{booktabs}
\usepackage{fancyhdr}
\usepackage[hidelinks]{hyperref}
\usetikzlibrary{arrows.meta,positioning,calc,shapes.geometric}
\usepackage[a4paper,margin=2.1cm,headheight=26pt,headsep=8pt]{geometry}

\definecolor{primary}{RGB}{146,95,160}
\definecolor{primarydark}{RGB}{92,52,104}
\definecolor{corcol}{RGB}{0,114,178}
\definecolor{astucecol}{RGB}{198,93,9}
\definecolor{bondcol}{RGB}{60,60,70}
\definecolor{piecol}{RGB}{198,93,9}

\tcbset{boxbase/.style={enhanced,breakable,boxrule=0.9pt,arc=2.5pt,
   left=3mm,right=3mm,top=2.2mm,bottom=2.2mm,
   fonttitle=\bfseries,coltitle=white,toptitle=0.6mm,bottomtitle=0.6mm}}
\newtcolorbox[auto counter]{cor}[1][]{boxbase,colback=corcol!5,colframe=corcol,
   title={\textsc{Corrigé}~\thetcbcounter\ifx\relax#1\relax\relax\else\textmd{ —\itshape\ #1}\fi}}
\newtcolorbox{astucebox}[1][]{boxbase,colback=astucecol!8,colframe=astucecol,
   title={\textsc{Astuce}\ifx\relax#1\relax\relax\else\textmd{ : \itshape #1}\fi}}

\pagestyle{fancy}\fancyhf{}
\fancyhead[L]{\small\textcolor{primarydark}{\textbf{Fiche 6 · Thème 5 — La résonance}}}
\fancyhead[R]{\small\textcolor{primarydark}{\textsc{Difficile · Corrigé}}}
\fancyfoot[C]{\small\thepage}
\renewcommand{\headrulewidth}{0.8pt}
\renewcommand{\headrule}{\hbox to\headwidth{\color{primary}\leaders\hrule height \headrulewidth\hfill}}
\setlength{\parindent}{0pt}\setlength{\parskip}{3pt}

\newcommand{\rep}{\textbf{\textcolor{corcol}{Réponse : }}}

\begin{document}
\begin{center}
{\Large\bfseries\color{primarydark}Thème 5 — Corrigé du niveau Difficile}
\end{center}
\medskip

\begin{cor}[l'ion carbonate \ce{CO3^2-}, pas à pas]
\begin{enumerate}[label=\alph*),leftmargin=2em,itemsep=3pt]
  \item \rep une structure de Lewis unique impose une double liaison \ce{C=O} sur \emph{un} oxygène
        précis et deux liaisons simples \ce{C-O^-} sur les deux autres. Elle rend donc les oxygènes
        \emph{différents}, alors que l'expérience montre qu'ils sont tous équivalents (liaisons de
        même longueur). Une seule formule ne peut pas rendre compte de cette symétrie : il faut
        plusieurs formes limites.
  \item \rep les \textbf{trois} formes limites se déduisent en plaçant la double liaison
        successivement sur chacun des trois oxygènes ; les deux oxygènes restants portent alors
        chacun une charge $\ominus$ (charge totale $-2$) :

\begin{center}
\begin{tikzpicture}[scale=0.7,every node/.style={font=\small}]
  \tikzset{
    Oatom/.style={font=\bfseries},
    Catom/.style={font=\bfseries},
    sbond/.style={bondcol,line width=1.1pt},
    dbond/.style={bondcol,double,double distance=1.7pt,line width=0.6pt},
  }
  % ---- Form 1 : double en haut ----
  \begin{scope}[xshift=0cm]
    \coordinate (C) at (0,0);
    \coordinate (T) at (90:1.3);
    \coordinate (BL) at (210:1.3);
    \coordinate (BR) at (330:1.3);
    \draw[dbond] (C)--(T);
    \draw[sbond] (C)--(BL);
    \draw[sbond] (C)--(BR);
    \node[Catom] at (C){C};
    \node[Oatom] at (T){O};\node[Oatom] at (BL){O};\node[Oatom] at (BR){O};
    \node[piecol] at ($(BL)+(-0.42,-0.12)$){$\ominus$};
    \node[piecol] at ($(BR)+(0.42,-0.12)$){$\ominus$};
  \end{scope}
  \node at (2.5,0){$\longleftrightarrow$};
  % ---- Form 2 : double en bas-gauche ----
  \begin{scope}[xshift=5cm]
    \coordinate (C) at (0,0);
    \coordinate (T) at (90:1.3);
    \coordinate (BL) at (210:1.3);
    \coordinate (BR) at (330:1.3);
    \draw[dbond] (C)--(BL);
    \draw[sbond] (C)--(T);
    \draw[sbond] (C)--(BR);
    \node[Catom] at (C){C};
    \node[Oatom] at (T){O};\node[Oatom] at (BL){O};\node[Oatom] at (BR){O};
    \node[piecol] at ($(T)+(0,0.42)$){$\ominus$};
    \node[piecol] at ($(BR)+(0.42,-0.12)$){$\ominus$};
  \end{scope}
  \node at (7.5,0){$\longleftrightarrow$};
  % ---- Form 3 : double en bas-droite ----
  \begin{scope}[xshift=10cm]
    \coordinate (C) at (0,0);
    \coordinate (T) at (90:1.3);
    \coordinate (BL) at (210:1.3);
    \coordinate (BR) at (330:1.3);
    \draw[dbond] (C)--(BR);
    \draw[sbond] (C)--(T);
    \draw[sbond] (C)--(BL);
    \node[Catom] at (C){C};
    \node[Oatom] at (T){O};\node[Oatom] at (BL){O};\node[Oatom] at (BR){O};
    \node[piecol] at ($(T)+(0,0.42)$){$\ominus$};
    \node[piecol] at ($(BL)+(-0.42,-0.12)$){$\ominus$};
  \end{scope}
\end{tikzpicture}
\end{center}
        (D'une forme à la suivante, un doublet libre d'un oxygène chargé forme une double liaison
        vers \ce{C}, tandis que la double liaison \ce{C=O} redevient un doublet libre sur un autre
        oxygène : \emph{seuls les électrons se déplacent}.)
  \item \rep les trois liaisons \ce{C-O} sont \textbf{identiques}, de longueur \textbf{intermédiaire}
        entre celle d'une liaison simple \ce{C-O} (plus longue) et celle d'une liaison double
        \ce{C=O} (plus courte). Chaque liaison a un caractère double de $\tfrac{1}{3}$.
  \item \rep par symétrie, la charge $-2$ est répartie également : chaque oxygène porte en moyenne
        \[ \frac{-2}{3}\approx -0{,}67. \]
\end{enumerate}
\end{cor}

\begin{cor}[l'ion perchlorate \ce{ClO4-}, pas à pas]
\begin{enumerate}[label=\alph*),leftmargin=2em,itemsep=3pt]
  \item \rep il y a \textbf{4} formes limites équivalentes : la charge $-1$ (et la liaison simple qui
        l'accompagne) peut se placer sur \emph{chacun} des quatre oxygènes équivalents.
  \item \rep en passant d'une forme à la suivante : un \textbf{doublet libre} de l'oxygène qui portait
        la charge $-1$ se rabat pour former une \textbf{double liaison} \ce{Cl=O} ; simultanément,
        l'une des doubles liaisons $\pi$ portées par le chlore \textbf{redevient un doublet libre}
        sur un autre oxygène, qui prend alors la charge $-1$. Les noyaux ne bougent pas.
  \item \rep \textbf{non}, la géométrie tétraédrique n'est pas affectée. La résonance ne déplace que
        des électrons : les positions des atomes (donc les angles \ce{O-Cl-O} $\approx 109{,}5^\circ$)
        sont les mêmes dans toutes les formes limites. La géométrie est une propriété de l'hybride,
        pas d'une forme particulière.
  \item \rep la charge $-1$ n'est localisée sur \emph{aucun} oxygène en particulier : elle est
        \textbf{délocalisée également} sur les quatre oxygènes, soit une charge moyenne de
        $-\tfrac{1}{4}$ par oxygène. C'est ce qui rend les quatre \ce{Cl-O} rigoureusement identiques.
\end{enumerate}
\end{cor}

\begin{cor}[l'ozone \ce{O3}, molécule de l'atmosphère]
\begin{enumerate}[label=\alph*),leftmargin=2em,itemsep=3pt]
  \item \rep les deux formes limites : \quad \ce{O=O-O <-> O-O=O}.
  \item \rep la molécule réelle est l'\textbf{hybride} de ces deux formes. Elle reste \textbf{coudée}
        (l'atome central porte un doublet non liant ; la géométrie ne change pas avec la résonance).
        Ses deux liaisons \ce{O-O} sont \textbf{identiques}, de longueur intermédiaire entre une
        simple et une double liaison (caractère double $\tfrac{1}{2}$ chacune).
  \item \rep la flèche $\longleftrightarrow$ ne décrit \emph{pas} un mouvement : la molécule ne
        « bascule » pas d'une structure à l'autre. Il n'existe \textbf{qu'une seule} espèce, réelle
        et stable, dont les électrons $\pi$ sont délocalisés en permanence ; les deux dessins ne sont
        que deux \emph{approximations} imparfaites de cette unique molécule. Un basculement entre
        deux espèces qui s'interconvertissent, lui, se noterait $\rightleftharpoons$ (équilibre
        chimique) et mettrait en jeu \emph{deux} espèces distinctes — ce qui n'est pas le cas ici.
\end{enumerate}
\end{cor}

\begin{astucebox}[garder le cap]
Trois automatismes pour ne pas tomber dans les pièges : \textbf{(1)}~seuls les \emph{électrons}
bougent d'une forme à l'autre, jamais les atomes ; \textbf{(2)}~la \emph{géométrie} est celle de
l'hybride et ne change jamais ; \textbf{(3)}~$\longleftrightarrow$ (une espèce délocalisée) n'est
\emph{jamais} $\rightleftharpoons$ (deux espèces en équilibre).
\end{astucebox}

\end{document}
